\documentclass[10pt,a4paper]{amsart}
\usepackage[margin=19mm]{geometry}
\usepackage{amsmath,amssymb,mathtools}
\usepackage[scr=rsfs,cal=euler]{mathalfa}
\usepackage{xfrac}
\usepackage[unicode,hidelinks,bookmarks=false]{hyperref}
\newcommand{\ie}{i.e.\ }
\newcommand{\cf}{cf.\ }
\newcommand{\resp}{respectively\ }
\newcommand{\Subsection}{\subsection}
\providecommand{\nobreakdash}{\nobreak\hbox{-}}
\setlength{\emergencystretch}{2em}
\multlinegap0pt
\usepackage{amssymb}
\usepackage{mathtools}
\newtheorem{proposition}{Proposition}
\newcommand{\dd}{\,\mathrm{d}}
\title{A Bessel-zero obstruction to hyperbolic complete monotonicity of noncentral chi-square densities}
\author{Domingos S. P. Salazar}
\date{Source: arXiv:2606.22066v1 (CC BY 4.0)}
\begin{document}
\maketitle

\noindent\emph{Source and licence.} A Bessel-zero obstruction to hyperbolic complete monotonicity of noncentral chi-square densities. Version arXiv:2606.22066v1 (CC BY 4.0), distributed by arXiv under the Creative Commons Attribution 4.0 International licence, http://creativecommons.org/licenses/by/4.0/, as recorded in the arXiv licence metadata for this exact version and shown on the abstract page https://arxiv.org/abs/2606.22066v1. Full licence notice: \texttt{licenses/arxiv-2606.22066-CC-BY-4.0.txt}.\par\smallskip

\noindent\textit{Editorial scope.} Counted unit: the article's proposition prop:bell-obstruction (source0.tex lines 356--390). The article's complete original proof of it is delivered with it (source0.tex lines 392--474, one \texttt{proof} environment of 83 lines). No mathematics is written, repaired or re-derived: the statement and the proof are the article's own lines, in the article's own notation, and no retained line is substituted or re-flowed.  Retained uncounted as the source-local context the proof needs: lines 325--328.  Nothing was omitted from any retained span, so the package carries no omission record.  No retained line contains a citation, so the wrapper carries no bibliography and no bibliography entry.  No retained line contains a figure environment, an image-inclusion command or drawing code, so no image is delivered and the package declares no asset dependency.  Editorial formatting only, in addition: an AMS \texttt{amsart} preamble that restates the article's own declarations and macros used by the retained lines, namely the environments \texttt{proposition} on separate counters, together with the standard packages those lines need.  The retained environments print in this document as Proposition 1 in order of appearance, whereas the article prints the same items with the numbering of its own full document.  The complete native source of the article as distributed in the arXiv e-print for this exact version is bundled unchanged as \texttt{sources/203320/source0.tex}.
\begin{equation}
\label{eq:log-Fa}
\log F_a(z)=\sum_{m=1}^{\infty}q_m z^m
\end{equation}

\begin{proposition}[Leading Bell obstruction]
\label{prop:bell-obstruction}
Let \(a,b,\theta>0\), and let \(p=p_{a,b,\theta}\).  Define
\begin{equation}
H_u(w)=p(uv(w))p(u/v(w)),
\qquad w>2.
\end{equation}
Let
\begin{equation}
d_1=b\left(\frac{\theta}{a}-1\right),
\qquad
d_m=m!q_m(b\theta)^m
\quad(m\ge2),
\end{equation}
where the \(q_m\)'s are defined by \eqref{eq:log-Fa}.  Then, for each fixed \(n\ge1\),
\begin{equation}
\label{eq:bell-leading}
\frac{1}{H_u(w)}
\frac{\dd^n}{\dd w^n}H_u(w)
=
u^n\mathcal B_n(d_1,\dots,d_n)
+
O(u^{n+1})
\qquad
(u\downarrow0),
\end{equation}
locally uniformly for \(w\) in compact subintervals of \((2,\infty)\).

Consequently, if \(p_{a,b,\theta}\) is HCM, then
\begin{equation}
\label{eq:hcm-necessary-bell}
(-1)^n\mathcal B_n(d_1,\dots,d_n)\ge0
\qquad(n\ge1).
\end{equation}
\end{proposition}

\begin{proof}
Fix a compact interval \(K\subset(2,\infty)\) and an integer \(n\ge1\).  On \(K\), the functions \(v(w)\) and \(v(w)^{-1}\) are bounded.  Since \(F_a(0)=1\), choose \(\rho>0\) such that \(F_a\) is zero-free on \(|z|<\rho\), and choose the branch of \(\log F_a\) there with value \(0\) at the origin.  For all sufficiently small \(u>0\), the two arguments \(b\theta uv(w)\) and \(b\theta u/v(w)\) lie in a compact subdisk of \(|z|<\rho\), uniformly for \(w\in K\).  Hence the Taylor series for \(\log F_a\), differentiated after composition with \(v(w)\) and \(v(w)^{-1}\) up to order \(n\), converges uniformly on \(K\).  In particular, all remainders below are locally uniform in \(w\).

For \(w>2\), let \(v=v(w)\).  Up to factors independent of \(w\),
\begin{equation}
H_u(w)
=
e^{-bu(v+v^{-1})}
F_a(b\theta uv)F_a(b\theta u/v).
\end{equation}
Taking logarithms and using \eqref{eq:log-Fa}, for sufficiently small \(u\) we have, uniformly on \(K\),
\begin{equation}
\log H_u(w)
=
C_u-buw+
\sum_{m=1}^{\infty}q_m(b\theta u)^mP_m(w),
\end{equation}
where \(C_u\) is independent of \(w\).

The analyticity just noted also controls the differentiated tail.  To see this explicitly, put
\begin{equation}
R_K=\sup_{w\in K}\max\{v(w),v(w)^{-1}\}.
\end{equation}
For each \(0\le r\le n\), differentiating \(P_m(w)=v(w)^m+v(w)^{-m}\) \(r\) times gives
\begin{equation}
\sup_{w\in K}|P_m^{(r)}(w)|
\le C_{K,r}m^rR_K^m
\qquad(m\ge1).
\end{equation}
Choose \(u_0>0\) so that \(|b\theta|u_0R_K\) lies inside the disk of convergence of \(\log F_a\).  Then, uniformly for \(0<u<u_0\) and \(w\in K\), the \(r\)-th derivative of the tail
\begin{equation}
\sum_{m\ge r+1}q_m(b\theta u)^mP_m(w)
\end{equation}
is bounded by
\begin{equation}
C_{K,r}\sum_{m\ge r+1}|q_m|\,m^r(|b\theta|u)^mR_K^m
=O(u^{r+1}).
\end{equation}
Thus the contribution of the terms \(m\ge r+1\) to the \(r\)-th \(w\)-derivative is \(O(u^{r+1})\), locally uniformly on \(K\).

Since \(P_m\) has degree \(m\) and leading coefficient \(1\),
\begin{equation}
\frac{\dd^r}{\dd w^r}P_m(w)=0
\quad(m<r),
\end{equation}
and
\begin{equation}
\frac{\dd^r}{\dd w^r}P_r(w)=r!.
\end{equation}
Thus, for \(r=1\),
\begin{equation}
\frac{\dd}{\dd w}\log H_u(w)
=
\left(q_1b\theta-b\right)u+O(u^2)
=
b\left(\frac{\theta}{a}-1\right)u+O(u^2)
=
d_1u+O(u^2),
\end{equation}
and for \(2\le r\le n\),
\begin{equation}
\frac{\dd^r}{\dd w^r}\log H_u(w)
=
r!q_r(b\theta)^r u^r+O(u^{r+1})
=
d_ru^r+O(u^{r+1}).
\end{equation}
The standard exponential derivative formula gives
\begin{equation}
\frac{H_u^{(n)}(w)}{H_u(w)}
=
\mathcal B_n
\left(
(\log H_u)'(w),
(\log H_u)''(w),
\dots,
(\log H_u)^{(n)}(w)
\right).
\end{equation}
Substituting the previous asymptotics into the weighted homogeneity of the complete Bell polynomials gives \eqref{eq:bell-leading}.  Explicitly, every monomial in \(\mathcal B_n\) has weighted degree \(n\), where the variable in position \(r\) has weight \(r\); replacing \((\log H_u)^{(r)}\) by \(d_ru^r+O(u^{r+1})\) therefore gives the leading term \(u^n\mathcal B_n(d_1,\dots,d_n)\) and an \(O(u^{n+1})\) remainder, locally uniformly on \(K\).

If \(p\) is HCM, then \((-1)^nH_u^{(n)}(w)\ge0\) for every \(u>0\), \(w>2\), and \(n\ge1\).  Since \(H_u(w)>0\), \((-1)^nH_u^{(n)}(w)/H_u(w)\ge0\).  Dividing \eqref{eq:bell-leading} by \(u^n>0\) and then letting \(u\downarrow0\), with any fixed \(w>2\), gives \eqref{eq:hcm-necessary-bell}.
\end{proof}

\end{document}
